\chapter{Post-Quantum Migration of Bitcoin and Ethereum}
\label{ch:blockchains}
\chaptermark{Post-Quantum Migration of Blockchains}

\section{Why this chapter matters}
Blockchains are one of the hardest post-quantum migration problems in existence, for two reasons the previous chapter only touched. First, the funds are secured by \emph{signatures} whose public keys are, sooner or later, \emph{permanently public} on an immutable ledger -- so ``harvest now, decrypt later'' becomes a much sharper ``store now, steal later.'' Second, there is no administrator: changing the rules of Bitcoin or Ethereum is a governance problem settled by rough consensus and forks, not a software update pushed to servers. This chapter looks at how the two largest chains are actually approaching the transition, as of mid-2026. Everything here is a moving target: the concrete proposals are drafts, not deployed consensus, and the reader should treat dates and statuses as a snapshot.

\section{Learning goals}
After reading this chapter, you should be able to:

\begin{itemize}
  \item explain which blockchain primitives a quantum computer breaks (Shor) versus merely weakens (Grover), and the public-key-exposure principle;
  \item describe Bitcoin's exposure by output type and the current draft proposals (BIP-360, BIP-361);
  \item describe Ethereum's exposure and its two-front plan: an emergency hard fork and account-abstraction-based migration;
  \item weigh the central controversy -- whether to freeze quantum-vulnerable coins -- and the honest uncertainty in the timeline.
\end{itemize}

\section{The blockchain threat model}
The immediate migration pressure is on signatures and pairing-based public-key systems.
Shor's algorithm (Chapter~\ref{ch:quantum}) threatens Bitcoin's ECDSA and Schnorr,
Ethereum's account ECDSA and consensus BLS signatures, and the public-key assumptions
behind SNARK and KZG systems~\cite{quantumhorizon2026,ethpqroadmap}. Grover gives
hashes a square-root rather than exponential weakening, but hashes still require
property- and output-length-specific assessment: HASH160, for example, has only a
nominal $2^{80}$ quantum preimage margin.

The subtlety that governs everything is \textbf{when a public key becomes visible}.
P2PK and P2TR expose a public key or tweaked public key when a UTXO is created;
P2PKH and P2WPKH usually reveal it on first spend, and reuse leaves later outputs
exposed. Other scripts must be assessed by their actual path and witness.

\begin{figure}[H]
\centering
\begin{tikzpicture}[node distance=6mm and 12mm, every node/.style={font=\scriptsize}]
  \node[flownode, align=center] (funded) {funded, never-spent address\\[-2pt]\scriptsize only a \emph{hash} of the public key is on chain};
  \node[flownode, align=center, right=of funded] (spend) {first spend\\[-2pt]\scriptsize the signature reveals the public key};
  \node[keynode, align=center, right=of spend] (exposed) {public key now on chain\\[-2pt]\scriptsize Shor can derive the private key};
  \draw[flowarrow] (funded) -- (spend);
  \draw[flowarrow] (spend) -- (exposed);
  \node[pqcteal] at ([yshift=-7mm]funded.south) {Grover only: safe};
  \node[pqcorange] at ([yshift=-7mm]exposed.south) {Shor: exposed};
\end{tikzpicture}
\caption{The public-key-exposure lifecycle. While an address has only ever received funds, its public key is hidden behind a hash and only Grover applies. The first spend publishes the public key permanently, exposing the account to Shor's algorithm. This is why exposed coins are considered at risk indefinitely: the ledger never forgets, so an attacker with a future quantum computer needs no fresh interception -- the target is already recorded.}
\label{fig:pubkey-exposure}
\end{figure}

\section{Bitcoin}
\subsection{What is exposed}
Bitcoin's exposure is precisely a function of output type. Coins in \emph{pay-to-public-key} (P2PK, the earliest, Satoshi-era outputs), raw multisig, and Taproot key-path outputs carry a public key on chain permanently and are long-exposure vulnerable. Coins in the hash-wrapped types (P2PKH, P2SH, P2WPKH) stay protected \emph{until the address is reused or spent}, at which point the key is revealed~\cite{bip360,chaincodepq2025}. BIP-361 states that, as of 1~March 2026, over \textbf{34\% of all bitcoin have revealed a public key on chain}~\cite{bip361}; an independent Chaincode Labs analysis puts the exposed fraction in the range of roughly one-fifth to one-half of supply~\cite{chaincodepq2025}. (Older and widely repeated figures attributing a specific amount to Satoshi-era P2PK are best treated as estimates rather than established fact.)

\subsection{The draft proposals}
Two Bitcoin Improvement Proposals anchor the response, and it is worth stressing that both are \textbf{Draft and not activated} as of mid-2026:

\begin{itemize}
  \item \textbf{BIP-360, ``Pay-to-Merkle-Root'' (P2MR)}~\cite{bip360} -- formerly circulated as ``Pay-to-Quantum-Resistant-Hash'' (P2QRH) -- introduces a new output type, proposed as a backward-compatible soft fork. In its current form it is essentially Taproot with the key-path spend removed, so that no bare public key is ever committed; notably, the current draft \emph{defers} the choice of the actual post-quantum signature scheme (naming ML-DSA and SLH-DSA only as examples) to a future proposal, rather than mandating one.
  \item \textbf{BIP-361, ``Post-Quantum Migration and Legacy Signature Sunset''}~\cite{bip361} -- the migration-and-deadline layer. It proposes a phased sunset of ECDSA/Schnorr: a first phase (on the order of three years after activation) that stops sending funds to legacy address types, and a later flag-day (around five years) that restricts legacy signatures, so that only a holder who can prove ownership by a quantum-safe route can still move the coins.
\end{itemize}

A recurring safety mechanism is \emph{commit--reveal}: a holder first publishes a
binding commitment and only later reveals it. It is designed to mitigate front-running
under the stated confirmation-depth, reorganization, fee, timeout, and recovery
assumptions; it is not an unconditional guarantee.

\subsection{The freeze controversy}
BIP-361's sunset implies that coins whose owners never migrate -- and coins for which no ownership proof is possible, notably P2PK -- become effectively unspendable, i.e.\ \emph{frozen}. This is the sharpest open dispute in Bitcoin's roadmap. Proponents argue a pre-announced sunset is defensive: a quantum attacker would otherwise \emph{steal} those coins outright, whereas a freeze at least preserves them. Critics call any forced migration confiscatory, contrary to Bitcoin's property-rights ethos, and warn that freezing long-dormant coins (including Satoshi's) sets a dangerous precedent. This is a genuine, unresolved political disagreement, not a settled plan, and should be read as such.

\section{Ethereum}
\subsection{What is exposed}
Ethereum addresses reveal an effectively public key once an ECDSA transaction has been
sent. BLS validator signatures, pairing curves behind SNARK verifiers, and KZG
commitments all rely on public-key assumptions threatened by Shor~\cite{ethpqroadmap,quantumhorizon2026}.
A comparison of quantum cost for secp256k1 and BLS12-381 needs group-specific resource
estimates; the 381-bit base-field name alone is not such an estimate. KZG migration
must distinguish blob retention from commitment binding, proof verification, and
consensus dependencies; blob pruning alone does not make the risk mild.

\subsection{The two-front plan}
Ethereum's response has an emergency track and a structural track; nothing post-quantum has yet shipped to mainnet.

\paragraph{The quantum-emergency hard fork.}
Vitalik Buterin's proposal ``How to hard-fork to save most users' funds in a quantum emergency''~\cite{buterinquantum2024} is the contingency plan. On evidence of large-scale theft, the chain would revert to the last pre-theft block, freeze legacy ECDSA accounts, and let users recover funds through a new transaction type carrying a \emph{STARK proof} that the sender knows the hash preimage secret behind their address -- authenticating by knowledge of a hash preimage rather than by an ECDSA signature. Because most users' recovery secret was never exposed on chain, this protects the majority, at the cost of a wallet-software update. Buterin notes the infrastructure could in principle be built now.

\paragraph{Account abstraction as the migration vehicle.}
The structural plan is to let accounts choose their own signature scheme. \textbf{EIP-7702} (``Set Code for EOAs,'' \emph{Final}, shipped in the 2025 Pectra upgrade)~\cite{eip7702} lets an externally-owned account delegate to contract code, turning it into a smart account with arbitrary verification logic -- although 7702 authorizations are themselves secp256k1-signed, so it is a stepping stone, not itself a post-quantum migration. The dedicated vehicle is \textbf{EIP-8141} (``Frame Transaction,'' \emph{Draft}, 2026, co-authored by Buterin)~\cite{eip8141}, which explicitly describes native account abstraction as ``a native off-ramp \dots{} to post-quantum secure systems,'' supporting a slot for an arbitrary (custom, post-quantum) signature scheme alongside the classical ones. At the consensus layer, the plan is to replace BLS aggregation with a hash-based signature of the XMSS/Winternitz family (Chapter~\ref{ch:hashots}), aggregated inside a STARK-based zero-knowledge virtual machine, and to move KZG commitments to STARK- or lattice-based alternatives; the Ethereum Foundation's stated assessment targets first-layer post-quantum infrastructure around 2029~\cite{ethpqroadmap}.

\subsection{On-chain verification, revisited}
The structural migration path may introduce post-quantum signature verification on
chain; the emergency path instead proposes a STARK proof of seed/preimage knowledge.
These are different authentication mechanisms. Their costs are scheme-, parameter-,
implementation-, chain-, and gas-schedule-dependent, so a benchmark must state those
details, including code commit and measurement date, before comparing gas figures.

\begin{table}[H]
\centering
\small
\begin{tabular}{@{}p{2.9cm}p{4.7cm}p{5.6cm}@{}}
\toprule
 & \textbf{Bitcoin} & \textbf{Ethereum} \\
\midrule
Fund-securing signature & ECDSA / Schnorr (secp256k1) & ECDSA (accounts); BLS12-381 (consensus) \\
Exposed today & P2PK, reused P2PKH, Taproot key-path & any account that has sent a transaction \\
Migration vehicle & new PQ output type (BIP-360) & account abstraction (EIP-7702, EIP-8141) \\
Emergency plan & legacy-signature sunset (BIP-361) & emergency hard fork with STARK recovery \\
Fork style & soft fork & protocol upgrades / emergency hard fork \\
Status (mid-2026) & Draft BIPs, not activated & research; one AA piece shipped; PQ infra $\sim$2029 \\
\bottomrule
\end{tabular}
\caption{Bitcoin and Ethereum post-quantum migration at a glance, as of mid-2026. Both chains face the same exposure principle but respond with different mechanisms and governance styles, and in both cases the concrete post-quantum proposals are drafts rather than deployed consensus.}
\label{tab:chain-migration}
\end{table}

\section{Timelines and honest uncertainty}
No quantum computer capable of breaking secp256k1 or BLS12-381 exists as of mid-2026.
A 2026 arXiv preprint, not a consensus survey, gives one set of model-dependent
estimates for cryptographically relevant quantum computing~\cite{quantumhorizon2026}.
A 2025 Google analysis of RSA-2048 illustrates a downward resource-estimation trend,
but is RSA-specific and not a direct estimate for Bitcoin or Ethereum curves~\cite{gidney2025}.

\section{Pitfalls}

\paragraph{Pitfall 1: assuming ``my address is a hash, so I am safe.''}
True only until the first spend. Any account that has transacted has exposed its public key; on Ethereum that is essentially every active account, and on Bitcoin it includes every reused address.

\paragraph{Pitfall 2: treating draft proposals as decided.}
BIP-360, BIP-361, and EIP-8141 are drafts; the freeze question in particular is contested. Reporting them as Bitcoin's or Ethereum's ``plan'' overstates consensus.

\paragraph{Pitfall 3: quoting a single gas figure.}
On-chain verification costs move between implementations, versions, and precompile assumptions. A number without a source and date is not meaningful.

\section{Summary}
Blockchains sharpen the quantum threat and complicate the response:

\begin{itemize}
  \item the danger is to signatures (Shor), not hashing (Grover); an exposed public key on an immutable ledger is a permanent, ``store now, steal later'' liability;
  \item \textbf{Bitcoin} proposes a new quantum-hardened output type (BIP-360) and a contested legacy-signature sunset (BIP-361), both draft soft forks;
  \item \textbf{Ethereum} pairs an emergency hard fork with STARK-based hash-preimage recovery and an account-abstraction migration path (EIP-7702 shipped, EIP-8141 draft), targeting first-layer post-quantum infrastructure around 2029;
  \item the timeline is uncertain but the required lead time is long, which is why both communities are debating the transition years before any quantum computer can execute the attack.
\end{itemize}

This is where cryptographic theory meets the messiest realities of deployment -- immutable history, decentralized governance, and irreversible loss -- and it is a fitting place to close: the mathematics is ready, and the harder work of migrating the systems that depend on it has only begun.
